Este trabalho apresenta o desenvolvimento de um sistema web voltado à gestão e reserva de espaços físicos na universidade, como mesas, salas e auditórios, destinado aos integrantes da comunidade acadêmica. O sistema busca otimizar o uso dos ambientes institucionais, promovendo a organização e o compartilhamento eficiente dos recursos disponíveis.

Para garantir uma representação fiel dos espaços, a Coordenadoria de Projetos de Engenharia e Arquitetura (COPEA) colaborou com a cessão das plantas baixas oficiais, permitindo a criação de uma interface de visualização que reproduz com precisão as dimensões e a disposição real dos ambientes.

O projeto foi desenvolvido utilizando tecnologias web modernas, com foco na usabilidade, acessibilidade e integração entre usuários e infraestrutura universitária. Como resultado, obteve-se uma ferramenta intuitiva e funcional, capaz de contribuir para o aprimoramento da gestão de espaços e o fortalecimento da convivência colaborativa na comunidade acadêmica. \hfill \break

\providecommand{\keywords}[1]{\textbf{\textit{Palavras Chave---}} #1}
\noindent\keywords{Sistema de Reservas; Gestão de Espaços; Universidade; Visualização de Ambientes;}
